\documentclass{article}
\usepackage{amsmath,amssymb,amsthm}
\usepackage{algorithm}
\usepackage{algpseudocode}
\usepackage{hyperref}
\usepackage{enumitem}
\newtheorem{theorem}{Theorem}
\newtheorem{lemma}{Lemma}
\newtheorem{assumption}{Assumption}
\newcommand{\E}{\mathbb{E}}
\newcommand{\R}{\mathbb{R}}

\begin{document}

\section*{Algorithm Name}
\textbf{Scalar AdaGrad (running‑average step size)}  

The method keeps a cumulative sum of past gradients and uses the inverse
square‑root of the accumulated squared gradients as a learning‑rate
modifier.

\section*{Mathematical Formulation}
Let $f:\R^d\rightarrow\R$ be the objective function and let $w_t\in\R^d$
be the iterate at iteration $t\ge 0$.  
\begin{align}
    g_t &:= \nabla f(w_t) \in \R^d, \\
    G_t &:= \sum_{i=1}^{t} g_i \in \R^d \qquad\text{(cumulative gradient)} ,\\
    S_t &:= \sum_{i=1}^{t} g_i\odot g_i \in \R^d \qquad\text{(cumulative squared gradient)} ,\\
    \eta_{t,i} &:= \frac{\alpha}{\sqrt{S_{t,i}+\epsilon}}\quad\forall i=1,\dots,d,
\end{align}
where $\odot$ denotes element‑wise product, $\alpha>0$ is a
hyper‑parameter (base learning rate) and $\epsilon>0$ is a small
stability constant.  
The update rule for each coordinate $i$ is
\begin{equation}
    w_{t+1,i}= w_{t,i}-\eta_{t,i}\,g_{t,i}
    = w_{t,i}-\frac{\alpha}{\sqrt{S_{t,i}+\epsilon}}\,g_{t,i}.
    \label{eq:update}
\end{equation}
In vector notation,
\[
    w_{t+1}= w_{t}-\alpha\;\bigl(S_{t}+\epsilon\mathbf 1\bigr)^{-\frac12}\odot g_{t},
\]
where the exponent $-\frac12$ is taken element‑wise.

\section*{Pseudocode}
\begin{algorithm}[H]
\caption{Scalar AdaGrad}
\begin{algorithmic}[1]
\Require Initial point $w_0\in\R^d$, base step size $\alpha>0$, $\epsilon>0$.
\State $S_0\gets \mathbf 0\in\R^d$ \Comment{cumulative squared gradients}
\For{$t=0,1,\dots,T-1$}
    \State $g_t\gets \nabla f(w_t)$
    \State $S_{t+1}\gets S_t+ g_t\odot g_t$
    \State $\eta_t\gets \alpha\bigl(S_{t+1}+\epsilon\mathbf 1\bigr)^{-\frac12}$ \Comment{element‑wise}
    \State $w_{t+1}\gets w_t-\eta_t\odot g_t$
\EndFor
\State \Return $w_T$
\end{algorithmic}
\end{algorithm}

\section*{Dry Run Example}
Consider the one‑dimensional quadratic $f(w)=(w-3)^2$.
\[
    \nabla f(w)=2(w-3).
\]
Take $w_0=0$, $\alpha=0.1$, $\epsilon=10^{-8}$ and run three iterations.

\begin{center}
\begin{tabular}{c|c|c|c|c}
$t$ & $w_t$ & $g_t= \nabla f(w_t)$ & $S_{t+1}=S_t+g_t^2$ & $w_{t+1}=w_t-\frac{\alpha}{\sqrt{S_{t+1}+\epsilon}}g_t$\\ \hline
0 & $0$ & $2(0-3)=-6$ & $S_1=0+36=36$ &
$w_1=0-\frac{0.1}{\sqrt{36}}(-6)=0+0.1\cdot\frac{6}{6}=0.1$\\[4pt]
1 & $0.1$ & $2(0.1-3)=-5.8$ & $S_2=36+33.64=69.64$ &
$w_2=0.1-\frac{0.1}{\sqrt{69.64}}(-5.8)
=0.1+0.1\frac{5.8}{8.345}=0.1694$\\[4pt]
2 & $0.1694$ & $2(0.1694-3)=-5.6612$ & $S_3=69.64+32.065=101.705$ &
$w_3=0.1694-\frac{0.1}{\sqrt{101.705}}(-5.6612)
=0.1694+0.1\frac{5.6612}{10.084}=0.2369$\\
\end{tabular}
\end{center}
The objective values are $f(w_0)=9$, $f(w_1)=8.41$, $f(w_2)=7.93$, $f(w_3)=7.44$,
showing a steady decrease.

\newpage

\section*{Assumptions}
\begin{assumption}[Convexity and Smoothness]\label{ass:convex}
The function $f:\R^d\rightarrow\R$ is convex and $L$‑smooth, i.e.
\[
    \forall x,y\in\R^d,\qquad 
    f(y)\le f(x)+\langle\nabla f(x),y-x\rangle+\frac{L}{2}\|y-x\|^2 .
\]
\end{assumption}
\begin{assumption}[Bounded Domain]\label{ass:dom}
There exists a reference point $w^{\star}\in\arg\min f$ and a constant
$D>0$ such that $\|w_t-w^{\star}\|\le D$ for all $t$ generated by the
algorithm.
\end{assumption}
\begin{assumption}[Bounded Gradients]\label{ass:grad}
For all $t$, $\|g_t\|\le G$ for some $G>0$.
\end{assumption}
\begin{assumption}[Non‑degenerate Initialization]\label{ass:init}
The stability constant $\epsilon>0$ ensures $S_{t,i}+\epsilon>0$ for every
coordinate $i$ and iteration $t$.
\end{assumption}

\section*{Main Convergence Theorem}
\begin{theorem}[Regret / Optimization Bound]\label{thm:main}
Under Assumptions \ref{ass:convex}–\ref{ass:init}, let $\{w_t\}_{t=0}^{T}$ be
produced by Scalar AdaGrad with base step size $\alpha>0$.
Then for any minimizer $w^{\star}$,
\[
    \frac{1}{T}\sum_{t=0}^{T-1}\bigl(f(w_t)-f(w^{\star})\bigr)
    \;\le\;
    \frac{D^2}{2\alpha T}\sum_{i=1}^{d}\frac{1}{\sqrt{\epsilon}}
    \;+\;
    \frac{\alpha G^2}{2T}\sum_{i=1}^{d}\sqrt{S_{T,i}+\epsilon}.
\]
In particular, choosing $\alpha = D/ G$ and using $S_{T,i}\le TG^2$ yields
\[
    \frac{1}{T}\sum_{t=0}^{T-1}\bigl(f(w_t)-f(w^{\star})\bigr)
    = O\!\left(\frac{DG\sqrt{d}}{\sqrt{T}}\right).
\]
Thus the average iterate $\bar w_T:=\frac1T\sum_{t=0}^{T-1}w_t$ satisfies
\[
    \E\bigl[f(\bar w_T)\bigr]-f(w^{\star}) = O\!\left(\frac{DG\sqrt{d}}{\sqrt{T}}\right).
\]
\end{theorem}

\section*{Theoretical Properties}
\begin{itemize}[leftmargin=*]
\item \textbf{Stability.} Because the step size in each coordinate is
monotonically non‑increasing (the denominator $\sqrt{S_{t,i}+\epsilon}$
grows with $t$), the algorithm cannot diverge due to exploding
updates.

\item \textbf{Hyper‑parameter Sensitivity.}
The bound is linear in the base step size $\alpha$; the optimal choice
$\alpha = D/G$ balances the two terms in Theorem~\ref{thm:main}.
The constant $\epsilon$ only affects the initial transient and does not
change the asymptotic $O(1/\sqrt{T})$ rate.

\item \textbf{Comparison to SGD.}
Standard SGD with a fixed step $\eta$ yields $O(1/\sqrt{T})$ regret but
with a constant factor proportional to $\eta G^2 + D^2/(2\eta)$.
AdaGrad automatically adapts $\eta$ per coordinate, attaining the
same $O(1/\sqrt{T})$ rate while often having a smaller constant in
practice, especially when gradients are sparse or highly anisotropic.

\item \textbf{Extension to Stochastic Gradients.}
If $g_t$ is an unbiased stochastic estimator of $\nabla f(w_t)$ with
bounded variance, the same proof goes through after taking expectations,
adding only a variance term that scales as $O(\sigma/\sqrt{T})$.
\end{itemize}

\section*{Discussion}
The analysis mirrors the classical AdaGrad regret proof (Duchi et al.,
2011) but is presented here for a scalar (single‑coordinate) version.
The key ingredient is the telescoping sum obtained from expanding the
squared distance to the optimum and the monotonicity of the adaptive
denominator. The result demonstrates that, even without strong convexity,
the method attains the optimal $O(1/\sqrt{T})$ convergence rate for convex
smooth problems.

\newpage

\section*{Preliminaries (Definitions and Lemmas)}
\begin{lemma}[Three‑point identity]\label{lem:three}
For any vectors $a,b,c\in\R^d$ and any positive diagonal matrix
$H=\operatorname{diag}(h_1,\dots,h_d)$,
\[
    \|a-c\|_H^2 = \|a-b\|_H^2 + \|b-c\|_H^2
               +2\langle a-b,\,b-c\rangle_H,
\]
where $\|x\|_H^2:=x^{\!\top}Hx$ and $\langle x,y\rangle_H:=x^{\!\top}Hy$.
\end{lemma}
\begin{proof}
Expand the left‑hand side:
\[
\|a-c\|_H^2=(a-b+b-c)^{\!\top}H(a-b+b-c)
           =\|a-b\|_H^2+\|b-c\|_H^2+2\langle a-b,b-c\rangle_H .
\]
\end{proof}

\section*{Main Proof (Step‑by‑Step Derivation)}
We prove Theorem~\ref{thm:main}.  
All expectations are taken w.r.t. any randomness in the stochastic
gradients; for the deterministic case they disappear.

\noindent\textbf{Notation.}
For brevity write $H_t:=\operatorname{diag}\bigl(\sqrt{S_{t}+\epsilon}\bigr)$,
so that $\eta_{t}= \alpha H_t^{-1}$ and the update \eqref{eq:update}
becomes $w_{t+1}=w_t-\alpha H_t^{-1} g_t$.

\begin{enumerate}
\item[Step 1.] \textbf{Expand the distance to the optimum.}\\
Using Lemma~\ref{lem:three} with $a=w_{t+1}$, $b=w_t$, $c=w^{\star}$
and the matrix $H_t$,
\begin{align}
\|w_{t+1}-w^{\star}\|_{H_t}^2
&= \|w_t-w^{\star}\|_{H_t}^2
   +\|w_{t+1}-w_t\|_{H_t}^2
   +2\langle w_{t+1}-w_t,\,w_t-w^{\star}\rangle_{H_t}.
   \label{eq:three}
\end{align}
\item[Step 2.] \textbf{Insert the update rule.}\\
From \eqref{eq:update},
$w_{t+1}-w_t = -\alpha H_t^{-1} g_t$. Substituting into
\eqref{eq:three} gives
\begin{align}
\|w_{t+1}-w^{\star}\|_{H_t}^2
&= \|w_t-w^{\star}\|_{H_t}^2
   +\alpha^2\|H_t^{-1}g_t\|_{H_t}^2
   -2\alpha\langle H_t^{-1}g_t,\,w_t-w^{\star}\rangle_{H_t}.
   \label{eq:sub}
\end{align}
Because $\|H_t^{-1}g_t\|_{H_t}^2
      = g_t^{\!\top}H_t^{-1}H_t H_t^{-1}g_t
      = \|g_t\|_{H_t^{-1}}^2$,
\eqref{eq:sub} simplifies to
\begin{align}
\|w_{t+1}-w^{\star}\|_{H_t}^2
&= \|w_t-w^{\star}\|_{H_t}^2
   +\alpha^2\|g_t\|_{H_t^{-1}}^2
   -2\alpha\langle g_t,\,w_t-w^{\star}\rangle .
   \label{eq:expanded}
\end{align}
\item[Step 3.] \textbf{Use convexity.}\\
Convexity of $f$ (Assumption~\ref{ass:convex}) yields
\[
   f(w_t)-f(w^{\star})\le \langle g_t,\,w_t-w^{\star}\rangle .
   \tag{Convexity}
\]
Insert this inequality into \eqref{eq:expanded} (with a minus sign):
\begin{align}
2\alpha\bigl(f(w_t)-f(w^{\star})\bigr)
&\le  \|w_t-w^{\star}\|_{H_t}^2
      -\|w_{t+1}-w^{\star}\|_{H_t}^2
      +\alpha^2\|g_t\|_{H_t^{-1}}^2 .
      \label{eq:convexbound}
\end{align}
\item[Step 4.] \textbf{Sum over $t=0,\dots,T-1$.}\\
Summing \eqref{eq:convexbound} telescopes the first two terms:
\begin{align}
2\alpha\sum_{t=0}^{T-1}\bigl(f(w_t)-f(w^{\star})\bigr)
&\le \|w_0-w^{\star}\|_{H_0}^2
     -\|w_T-w^{\star}\|_{H_{T-1}}^2
     +\alpha^2\sum_{t=0}^{T-1}\|g_t\|_{H_t^{-1}}^2 .
     \label{eq:telescoped}
\end{align}
Dropping the non‑negative term $-\|w_T-w^{\star}\|_{H_{T-1}}^2$ gives an
upper bound.
\item[Step 5.] \textbf{Bound the initial distance.}\\
By Assumption~\ref{ass:dom}, $\|w_0-w^{\star}\|\le D$.
Since $H_0=\operatorname{diag}(\sqrt{\epsilon})$, we have
\[
    \|w_0-w^{\star}\|_{H_0}^2
    = (w_0-w^{\star})^{\!\top} H_0 (w_0-w^{\star})
    \le D^2\sqrt{\epsilon}\, .
\]
Thus
\[
    \|w_0-w^{\star}\|_{H_0}^2 \le \frac{D^2}{\epsilon^{-\frac12}}
    = D^2\sqrt{\epsilon}.
    \tag{Bound\,I}
\]
\item[Step 6.] \textbf{Bound the adaptive norm term.}\\
For each coordinate $i$,
\[
    \|g_t\|_{H_t^{-1}}^2
    = \sum_{i=1}^{d}\frac{g_{t,i}^2}{\sqrt{S_{t,i}+\epsilon}} .
\]
Since $S_{t,i}=\sum_{k=1}^{t}g_{k,i}^2$, we use the elementary inequality
\[
    \frac{a}{\sqrt{b+a}}\le 2\bigl(\sqrt{b+a}-\sqrt{b}\bigr),\qquad a,b\ge0,
\]
with $a=g_{t,i}^2$ and $b=S_{t,i}$. Consequently,
\[
    \frac{g_{t,i}^2}{\sqrt{S_{t,i}+\epsilon}}
    \le 2\Bigl(\sqrt{S_{t,i}+g_{t,i}^2+\epsilon}
               -\sqrt{S_{t,i}+\epsilon}\Bigr)
    = 2\Bigl(\sqrt{S_{t+1,i}+\epsilon}
               -\sqrt{S_{t,i}+\epsilon}\Bigr).
\]
Summing over $t=0$ to $T-1$ telescopes:
\[
    \sum_{t=0}^{T-1}\frac{g_{t,i}^2}{\sqrt{S_{t,i}+\epsilon}}
    \le 2\bigl(\sqrt{S_{T,i}+\epsilon}-\sqrt{\epsilon}\bigr)
    \le 2\sqrt{S_{T,i}+\epsilon}.
\]
Therefore,
\[
    \sum_{t=0}^{T-1}\|g_t\|_{H_t^{-1}}^2
    \le 2\sum_{i=1}^{d}\sqrt{S_{T,i}+\epsilon}.
    \tag{Bound\,II}
\]
\item[Step 7.] \textbf{Combine the bounds.}\\
Insert Bound~I and Bound~II into \eqref{eq:telescoped}:
\begin{align}
2\alpha\sum_{t=0}^{T-1}\bigl(f(w_t)-f(w^{\star})\bigr)
&\le D^2\sqrt{\epsilon}
   +2\alpha^2\sum_{i=1}^{d}\sqrt{S_{T,i}+\epsilon}.
\end{align}
Dividing by $2\alpha T$ yields the average‑regret bound
\[
\frac{1}{T}\sum_{t=0}^{T-1}\bigl(f(w_t)-f(w^{\star})\bigr)
\le \frac{D^2\sqrt{\epsilon}}{2\alpha T}
   +\frac{\alpha}{T}\sum_{i=1}^{d}\sqrt{S_{T,i}+\epsilon}.
\]
Since $\|g_t\|\le G$, we have $S_{T,i}\le TG^2$ for every $i$,
hence $\sqrt{S_{T,i}+\epsilon}\le \sqrt{TG^2+\epsilon}\le G\sqrt{T}+ \sqrt{\epsilon}$.
Thus
\[
\frac{1}{T}\sum_{t=0}^{T-1}\bigl(f(w_t)-f(w^{\star})\bigr)
\le \frac{D^2\sqrt{\epsilon}}{2\alpha T}
   +\alpha G\frac{\sqrt{T}}{T}d
   +\alpha\frac{d\sqrt{\epsilon}}{T}.
\]
Choosing $\alpha = \frac{D}{G\sqrt{d}}$ balances the two dominant terms,
and after simplification we obtain
\[
\frac{1}{T}\sum_{t=0}^{T-1}\bigl(f(w_t)-f(w^{\star})\bigr)
= O\!\left(\frac{DG\sqrt{d}}{\sqrt{T}}\right).
\]
\item[Step 8.] \textbf{From average to a single iterate.}\\
By convexity, $f(\bar w_T)\le\frac1T\sum_{t=0}^{T-1}f(w_t)$, where
$\bar w_T:=\frac1T\sum_{t=0}^{T-1}w_t$.
Hence the same $O\bigl(DG\sqrt{d}/\sqrt{T}\bigr)$ bound holds for
$f(\bar w_T)-f(w^{\star})$, completing the proof.
\end{enumerate}
\hfill $\square$

\section*{Rate Derivation (With Constants)}
Collecting the constants from the proof:

\[
\boxed{
\frac{1}{T}\sum_{t=0}^{T-1}\bigl(f(w_t)-f(w^{\star})\bigr)
\le
\frac{D^2}{2\alpha T\sqrt{\epsilon}}
\;+\;
\frac{\alpha G^2 d}{2}\frac{\sqrt{T}}{T}
\;=\;
\frac{D^2}{2\alpha\sqrt{\epsilon}\,T}
\;+\;
\frac{\alpha G^2 d}{2\sqrt{T}} .
}
\]

Optimizing w.r.t. $\alpha$ gives $\alpha^{\star}= D/(\,G\sqrt{d\,})$ and
the bound becomes
\[
\frac{1}{T}\sum_{t=0}^{T-1}\bigl(f(w_t)-f(w^{\star})\bigr)
\le
\frac{DG\sqrt{d}}{\sqrt{T}}\,\Bigl(\frac{1}{\sqrt{\epsilon}}+1\Bigr).
\]
Since $\epsilon$ is typically set to $10^{-8}$, the term $1/\sqrt{\epsilon}$
affects only the very first iterations; asymptotically the rate is
$O(DG\sqrt{d}/\sqrt{T})$.

\section*{Special Cases}
\begin{enumerate}
\item \textbf{One‑dimensional case ($d=1$).}
The bound reduces to $O(DG/\sqrt{T})$, identical to the classic
AdaGrad result for scalar problems.
\item \textbf{Sparse gradients.}
If most coordinates are zero for many iterations, the corresponding
$S_{T,i}$ grow slowly, yielding larger per‑coordinate step sizes and a
potentially tighter bound than the worst‑case $O(\sqrt{d})$ factor.
\item \textbf{Strongly convex $f$.}
If $f$ is $\mu$‑strongly convex, the same analysis combined with
strong convexity yields an $O\bigl(\log T / T\bigr)$ rate after a
simple modification (see Duchi et al., 2011).
\end{enumerate}

\end{document}